\section{Arithmetic selection before complex nonresonance}\label{sec:selection}
The order of choices is essential. We first choose the finite list $\F$
from Proposition~\ref{prop:all-genus} at one suitable complex parameter
triple. We then choose integral $B,Q$ by finitely many open conditions.
Only after their choice do we impose nonresonance at a complex $A$.

\begin{lemma}\label{lem:finite-open}
For the fixed finite list $\F$, there is a nonempty open set
$\mathcal U\subset T\times\SL_r$ of parameters $(B,Q)$ with the following
properties:
\begin{enumerate}[label=(\roman*)]
\item On a nonempty open subset of the $A$-torus, $\rho_A$ is a smooth
irreducible character point and the selected differentials span its
cotangent space.
\item For each band curve in $\F$, its characteristic-coefficient map
from the $A$-torus to $T/W$ is dominant, where $W$ is the symmetric group.
\end{enumerate}
\end{lemma}
\begin{proof}
Put $q=r^2-1$. Use a fixed basis of the Lie algebra to trivialize the
tangent bundles of the matrix factors. At the initial suitable point, the
linear map consisting of the relation derivative together with the
selected invariant derivatives has kernel the $q$-dimensional conjugation
space. Its rank is therefore $2gq-q$. A corresponding minor is nonzero.
All entries are regular functions of $(A,B,Q)$: inverses are regular on
$\SL_r$, and all words are finite. If quotient charts were used instead,
one could clear their nonzero denominators with the same conclusion.
Include minors or products expressing relation submersion and the explicit
irreducibility conditions. This produces a nonzero regular function on
$T\times T\times\SL_r$ that guarantees (i). Expand it as a Laurent
polynomial in $A$. The locus of $(B,Q)$ where at least one coefficient
is nonzero is a nonempty open; there it is nonzero on a nonempty open
of the $A$-torus.

Along \eqref{eq:family}, each band in $\F$ has the same characteristic
polynomial as $A^kBC^l$ for some $k,l\geq1$. This includes $u_0$ by
cyclic rotation and commutation of $A$ with $B$; for the extra-handle
curves the extra matrix is the identity. At $Q=I$ the map in question is
\[
 T\longrightarrow T\longrightarrow T/W,\qquad
 A\longmapsto A^kBC_\mu^l\longmapsto(e_1,\ldots,e_{r-1}).
\]
The first arrow is finite surjective and the second is the finite Weyl
quotient. In characteristic zero its differential has full rank somewhere.
One such Jacobian minor is a Laurent polynomial in $A$, regular in $B,Q$.
The condition that it is not identically zero in $A$ is nonempty open
in $(B,Q)$ and implies dominance. This witness keeps exactly the fixed
eigenvalues $\mu_i$; no change of conjugacy class is being made.
There are only finitely many curves, so their opens and the open in (i)
intersect in the irreducible parameter space $T\times\SL_r$.
\end{proof}

Choose $(B,Q)\in\mathcal U$ with $B=B(m)$ and $Q\in\SL_r(\Z)$ by
Lemma~\ref{lem:arithmetic-dense}, and fix $C=QC_\mu Q^{-1}$. Then
\begin{equation}\label{eq:seed}
 \rho_0=(I,B,C,I,I,\ldots,I)
\end{equation}
is an actual representation into $\SL_r(\OO_L)$: every commutator is
identity. Let
\[
 Z=\overline{\Mod(\Sigma_g)\cdot[\rho_0]}^{\mathrm{Zar}}\subset\X.
\]
The twist $a\mapsto ab$ has powers sending $\rho_0$ to
$(B^n,B,C,I,\ldots,I)$. Since $B$ generates $T$ densely, continuity of
the quotient morphism gives
\begin{equation}\label{eq:torus-in-closure}
 [\rho_A]\in Z\qquad\text{for every }A\in T.
\end{equation}

\begin{lemma}\label{lem:choose-x}
There is an $A\in T$ such that $x=[\rho_A]$ is a smooth irreducible
character point where the selected differentials span the cotangent
space, and every curve in $\F$ has regular semisimple monodromy whose
powers are Zariski dense in its centralizer torus.
\end{lemma}
\begin{proof}
In a diagonal torus, failure of cyclic density is equivalent to
$\chi(t)=1$ for some nontrivial character
$\chi\in\Z^r/\Z(1,\ldots,1)$. For each such character its kernel is a
proper closed subset of $T$; the finite Weyl image of the union of its
conjugates is a proper closed subset of $T/W$. Properness follows from
dimension, since the quotient is finite. There are countably many such
characters. Also remove the discriminant locus in $T/W$ explicitly.
For each band map, dominance from Lemma~\ref{lem:finite-open} makes all
these inverse images proper closed subsets of the $A$-torus. For the
curve $a$ the spectral map of $A$ itself is dominant. For the curve $c$,
regular semisimplicity and nonresonance already follow from
Lemma~\ref{lem:units}, independently of $A$.

Intersect the nonempty open from Lemma~\ref{lem:finite-open}(i) with the
complements of this countable collection of proper closed sets, for all
curves in the finite list. Such an intersection is nonempty by the Baire
theorem: on any small closed ball contained in the open smooth torus,
proper complex algebraic subsets are closed with empty interior. The
selected $A$ has every required property, including avoidance of the
discriminant for each band.
\end{proof}

The $A$ of Lemma~\ref{lem:choose-x} is not asserted to be arithmetic.
This is harmless, because \eqref{eq:torus-in-closure} places it inside
the closure of the orbit of the single integral representation
\eqref{eq:seed}. Neither an arithmetic Baire argument nor a simultaneous
Hilbert-irreducibility assertion is used.
